% Part: computability
% Chapter: computability-theory
% Section: normal-form

\documentclass[../../../include/open-logic-section]{subfiles}

\begin{document}

% SEGMENT OLP-0231-S01

\olfileid{cmp}{thy}{nfm}
\olsection{இயல்புவடிவத் தேற்றம்}

கணிக்கத்தக்க சார்புகளின் வரையறைகளை நம்மால் விவரிக்கவும், ஏதேனும்
உள்ளீட்டில் அச்சார்பின் கணக்கீட்டின் முழுப் பதிவுக்கான சாத்தியமான
விளக்கம் சரியா என்று சோதிக்கவும் முடிகிறது என்று கொள்க. அப்போது
கணக்கீட்டு மாதிரியைச் சாராமல், எந்த உள்ளீடு~$x$ இலுள்ள எந்த
கணிக்கத்தக்க சார்பு~$f$ இன் மதிப்பையும் பின்வருமாறு தீர்மானிக்கலாம்
என்பது நியாயமானது:
\begin{enumerate}
  \item கணக்கீட்டுப் பதிவுகளின் சாத்தியமான அனைத்து விளக்கங்களிலும்
  தேடுக.
  \item கொடுக்கப்பட்ட பதிவு~$f(x)$ இன் கணக்கீட்டுப் பதிவா என்று
  சோதிக்க.
  \item சரியான பதிவைக் கண்டிருந்தால் அதிலிருந்து~$f(x)$ இன்
  மதிப்பைப் பிரித்தெடுக்க.
\end{enumerate}
இது உண்மையில் மெய் என்பதே க்ளீனியின் இயல்புவடிவத் தேற்றத்தின்
உள்ளடக்கம்.

\begin{thm}[க்ளீனியின் இயல்புவடிவத் தேற்றம்]
\ollabel{thm:normal-form}
பின்வரும் பண்புடைய முதன்மை மீள்வரையறை உறவு~$T(e, x, s)$ ஒன்றும்
முதன்மை மீள்வரையறைச் சார்பு~$U(s)$ ஒன்றும் உள்ளன: $f$ ஏதேனும்
பகுதியாகக் கணிக்கத்தக்க சார்பு எனில், ஏதேனும்~$e$ க்கு, ஒவ்வொரு~$x$
க்கும்
\[
f(x) \simeq U(\umin{s}{T(e, x, s)})
\]
பொருந்தும்.
\end{thm}

\begin{proof}[நிரூபிப்பு வரைவு]
எந்தக் கணக்கீட்டு மாதிரிக்கும், கணிக்கத்தக்க சார்பு~$f$ இன் விளக்கத்தைத்
துல்லியமாக வரையறுத்து, அத்தகைய விளக்கத்தை இயல் எண்~$e$ கொண்டு
குறியீடாக்கலாம். சுட்டெண்~$e$ உடைய சார்பை உள்ளீடு~$x$ இல்
கணிக்கும் செயல்முறையைப் பதிவுசெய்யும் “கணக்கீட்டுத் தொடர்”
என்ற கருத்தையும் துல்லியமாக வரையறுக்கலாம். அத்தகைய கணக்கீட்டுத்
தொடரையும் ஓர் எண்~$s$ ஆகக் குறியீடாக்கலாம். பின்வருவன அமையுமாறு
இதைச் செய்யலாம்:
\begin{enumerate}
  \item சுட்டெண்~$e$ உடைய சார்பின் உள்ளீடு~$x$ க்கான கணக்கீட்டுத்
  தொடரை ஓர் எண்~$s$ குறியீடாக்கினால், எனில் மட்டுமே, பொருந்தும்
  உறவு $T(e, x, s)$; மேலும்
  \item குறியீடு~$s$ உடைய கணக்கீட்டுத் தொடரை அக்கணக்கீட்டின்
  இறுதி விளைவுக்குக் கொண்டுசெல்லும் சார்பு $U(s)$
\end{enumerate}
ஆகிய இரண்டும் கணிக்கத்தக்கவை. உண்மையில் உறவு~$T$ மற்றும்
சார்பு~$U$ ஆகியவை முதன்மை மீள்வரையறைக்குரியவை.
\end{proof}

\begin{explain}
இயல்புவடிவத் தேற்றத்தின் துல்லியமான நிரூபிப்பைத் தர, ஒரு கணக்கீட்டு
மாதிரியை நிலைப்படுத்தி, கணிக்கத்தக்க சார்புகளின் விளக்கங்களையும்
கணக்கீட்டுத் தொடர்களையும் விரிவாகக் குறியீடாக்கி, உறவு~$T$ மற்றும்
சார்பு~$U$ முதன்மை மீள்வரையறைக்குரியவை என்று சரிபார்க்க வேண்டும்.
பெரும்பாலான பயன்பாடுகளுக்கு, $T$ மற்றும்~$U$ கணிக்கத்தக்கவை என்றும்
$U$ முழுச் சார்பு என்றும் இருந்தால் போதும்.

இயல்புவடிவத் தேற்றத்தின் நிரூபிப்பைப் பின்வரும் முழக்க வடிவில்
நினைவில் கொள்வது சிறந்தது: $\umin{s}{T(e, x, s)}$ ஆனது
சுட்டெண்~$e$ உடைய சார்பின் உள்ளீடு~$x$ க்கான கணக்கீட்டுத் தொடரைத்
தேடுகிறது; அத்தொடரைக் கண்டால் கணக்கீட்டுத் தொடரின் வெளியீட்டை
$U$ திருப்பித்தருகிறது.
\end{explain}

கணக்கீட்டு மாதிரி பகுதி மீள்வரையறைச் சார்புகளாக இருந்தால்
(க்ளீனி முதலில் பயன்படுத்தியது இதுவே), எந்தச் சார்பின் வரையறைக்கும்
வரம்பற்ற தேடலின் ஒரே ஒரு பயன்பாடு, அதாவது ஒரே ஒரு
$\umin{y}{f(\vec x, y)}$ செயற்குறி மட்டுமே தேவை என்பதை இது
காட்டுகிறது. இப்பொருளில், எந்த பகுதி மீள்வரையறைச் சார்புக்கும் ஓர்
இயல்புவடிவம் உள்ளது என்று இது காட்டுகிறது.

$T$ மற்றும் $U$ ஐப் பயன்படுத்தி $\cfind{0}$, $\cfind{1}$,
$\cfind{2}$, \dots என்ற எண்ணிடலை வரையறுக்கலாம். இனி,
$T$ மற்றும்~$U$ க்கு ஏற்ற ஒரு தேர்வை நிலைப்படுத்தியுள்ளோம் என்று
கொண்டு, பின்வரும் சமன்பாட்டை $\cfind{e}$ இன்
\emph{வரையறையாகக்} கொள்வோம்:
\[
\cfind{e}(x) \simeq U(\umin{s}{T(e,x,s)})
\]

மற்றொரு பயனுள்ள உண்மை இதோ:

\begin{thm}
ஒவ்வொரு பகுதியாகக் கணிக்கத்தக்க சார்பும் முடிவிலி எண்ணிக்கையிலான
சுட்டெண்களைக் கொண்டுள்ளது.
\end{thm}

மீண்டும், இது உள்ளுணர்வாகத் தெளிவு. ஒரு கணிக்கத்தக்க சார்பின்
(விளக்கம்) எதுவும் தரப்பட்டால், அதே சார்பைக் (அதே உள்ளீடு--வெளியீட்டுச்
சோடியை) கணிக்கும் வேறு விளக்கத்தை உருவாக்கலாம்; எடுத்துக்காட்டாக,
கணக்கீட்டின்மீது எந்த விளைவுமில்லாத ஒன்றை முதலில் செய்யலாம்
($0 = 0$ எனச் சோதித்தல், அல்லது $5$ வரை எண்ணுதல் முதலியவை).
மாற்றப்பட்ட விளக்கத்தின் சுட்டெண் மூலச் சுட்டெண்ணிலிருந்து எப்போதும்
வேறுபடும். இரண்டும் சற்றே வேறுபட்ட முறையில் கணிக்கப்பட்ட அதே சார்பின்
சுட்டெண்கள்.

\end{document}
